医学影像模型需要随数据来源、类别、器官和成像模态的变化持续更新，同时保持已有分析能力。这一过程受到精细标注获取困难、历史样本访问受限和跨中心数据难以集中等条件的制约。本文设计 MedCL 医学影像持续学习评测平台，建立持续分割基准，并研究跨中心弱监督分割、类别扩展下的无回放联邦分类和跨器官、跨模态持续配准。主要工作如下。

（1）设计 MedCL 评测平台，并构建持续医学图像分割基准。平台围绕任务顺序、阶段模型、固定测试集和评分规则组织分割、分类与配准评价，以阶段--任务性能矩阵记录模型表现，并为联邦场景设计分客户端评价。平台部分包括需求分析、系统架构、评测接口和测试方案；本文算法实验独立于平台完成。分割基准覆盖六中心域变化、心脏结构类别扩展和跨器官任务变化，分别评价最终性能、历史保持、当前学习、前向泛化与资源开销。六中心实验中，经验回放的最终平均 Dice 为 $0.750\pm0.021$，高于弹性权重巩固的 $0.672\pm0.021$，而二者的扩展前向迁移分别为 $0.256\pm0.150$ 和 $0.255\pm0.152$，表明最终性能的提高未必伴随前向泛化的改善。

（2）提出结合监督增强与特征一致性约束的弱监督持续分割方法 ZSDERpp。在静态弱监督方法 ZScribbleSeg 的基础上，利用全局一致性和空间先验补充稀疏涂鸦之外的监督约束，并结合有限历史图像--涂鸦回放与参考特征保持，缓解跨中心顺序更新中的遗忘。实验采用固定模拟涂鸦训练，使用密集验证标签选择模型，并以密集测试标签评价。六中心序列中，ZSDERpp 的最终平均 Dice 为 $0.701$，高于增强涂鸦顺序训练的 $0.510$；后向迁移率由 $-0.324$ 改善为 $-0.152$，但前向泛化未取得最高值。

（3）提出无回放联邦持续分类方法 FedSubMerge 及其自适应版本\linebreak \mbox{FedSubMerge-AD}。FedSubMerge 在服务器端融合客户端主梯度子空间，通过正交补投影约束后续共享参数更新；FedSubMerge-AD 按网络层选择相关客户端，以减轻不相关历史方向对当前学习的限制。实验采用类别批次持续到达、测试任务身份已知的多头分类设置。PathMNIST 的 Dirichlet 参数 $\alpha=0.3$ 时，FedSubMerge-AD 的最终平均准确率为 $93.42\%$，较无回放比较方法 FOT 提高 $7.22$ 个百分点；FedSubMerge 的后向迁移为 $-0.86$ 个百分点，优于 FOT 的 $-6.30$ 个百分点。完全参与时，阶段末子空间的额外通信量约为 20 个常规模型通信轮次传输量的 $3.0\%$--$3.5\%$。方法保留模型及子空间摘要，不回放历史原始样本；子空间交换本身不提供形式化隐私保证。

（4）提出结合有限经验回放、一阶元持续学习和锐度感知优化的持续配准方法 SAMCL。该方法以少量历史图像对参与后续训练，在器官、模态和配准关系变化时兼顾新任务适应与旧任务保持。实验依次学习脑部 MR--MR、腹部 CT--CT、肺部 CT--CT 和腹部 MR--CT 配准。与直接顺序训练相比，SAMCL 将脑部任务的最终 Dice 由 $0.348$ 提高到 $0.697$，后向变化由 $-0.395$ 改善为 $-0.037$；腹部 CT--CT 的后向变化由 $-0.129$ 改善为 $-0.039$。SAMCL 与元经验回放在不同任务上各有优势。

上述研究表明，评价医学影像持续学习既要考察最终性能，也要区分当前任务学习、历史保持和前向泛化。本文在保留各任务数据、模型与指标差异的基础上，给出评测平台设计及三类任务的实验研究，为训练信息受限条件下的医学影像模型持续更新提供参考。